% sections/appendix.tex
% Included from main.tex via % sections/appendix.tex
% Included from main.tex via % sections/appendix.tex
% Included from main.tex via % sections/appendix.tex
% Included from main.tex via \input{sections/appendix}

\appendix
\section{Appendix}
\subsection{System Configuration}
\label{app:config}


\begin{figure}[!htbp]
\centering
\includegraphics[width=\columnwidth, height=0.3\textheight, keepaspectratio]{figures/asset.jpg}
\caption{%
  Custom assets imported into CARLA-Air through the extensible asset pipeline.
  \textbf{Top:} a four-wheeled mobile robot with onboard LiDAR, imported from
  an external FBX model.
  \textbf{Bottom:} a custom electric sport car with user-defined vehicle
  dynamics.
  Both assets operate within the shared simulation world alongside all
  built-in CARLA traffic and AirSim aerial agents, and are visible to all
  sensor modalities.%
}
\label{fig:custom_assets}
\end{figure}


\paragraph{Reference hardware.}
All experiments in this report were conducted on the following configuration:
Ubuntu 20.04/22.04~LTS, NVIDIA RTX~A4000 (16\,GB GDDR6), AMD Ryzen~7 5800X
(8-core, 4.7\,GHz), 32\,GB DDR4-3200.

\paragraph{Software stack.}
CARLA~0.9.16, AirSim~1.8.1 (final stable open-source release),
Unreal Engine~4.26, Python~3.8+.

\paragraph{Network configuration.}
Table~\ref{tab:ports} lists the default port assignments.
Both RPC servers bind to \texttt{localhost} by default; remote connections
require explicit IP configuration.

\begin{table}[!htbp]
\centering
\caption{Default network port assignments.}
\label{tab:ports}
\footnotesize
\setlength{\tabcolsep}{4pt}
\renewcommand{\arraystretch}{1.08}
\begin{tabular}{@{}llr@{}}
\toprule
\textbf{Service} & \textbf{Protocol} & \textbf{Port} \\
\midrule
CARLA RPC Server   & TCP & 2000  \\
CARLA Streaming    & UDP & 2001  \\
AirSim RPC Server  & TCP & 41451 \\
\bottomrule
\end{tabular}
\end{table}

\paragraph{Distribution.}
The prebuilt binary package is approximately 19\,GB and includes a one-command
launcher (\texttt{CarlaAir.sh}).
The source distribution is approximately 651\,MB and is released under the MIT
license.


\subsubsection{API Compatibility Summary}
\label{app:api}

Table~\ref{tab:api_compat} summarizes the API compatibility status and test
coverage of the current CARLA-Air release.
All 89 automated CARLA API tests pass without modification; the full AirSim
flight control and sensor access API has been verified through manual and
scripted testing.
A total of 63 ROS\,2 topics are published across both simulation backends.

\begin{table}[!htbp]
\centering
\caption{API compatibility and test coverage.}
\label{tab:api_compat}
\small
\setlength{\tabcolsep}{3pt}
\renewcommand{\arraystretch}{1.08}
\resizebox{\columnwidth}{!}{%
\begin{tabular}{@{}lp{4.5cm}@{}}
\toprule
\textbf{Component} & \textbf{Status} \\
\midrule
CARLA API           & 89/89 automated tests passing \\
AirSim API          & Full flight control and sensor access verified \\
ROS\,2 topics       & 63 total (43 CARLA + 14 AirSim + 6 generic) \\
\midrule
\multicolumn{2}{@{}l}{\textit{Key upstream scripts confirmed working}} \\
\texttt{manual\_control.py}    & CARLA manual driving interface \\
\texttt{automatic\_control.py} & CARLA autopilot demonstration \\
\texttt{dynamic\_weather.py}   & CARLA weather preset cycling \\
\texttt{hello\_drone.py}       & AirSim basic flight demonstration \\
\bottomrule
\end{tabular}}%
\end{table}


\subsection{Upstream Source Modifications}
\label{app:source_mods}

CARLA-Air is designed to minimize modifications to the upstream CARLA codebase.
The integration touches only two header files and one source file, totaling
approximately 35 lines of changes.
All other integration code is purely additive, residing within the aerial
simulation plugin as the \texttt{CARLAAirGameMode} class
(${\sim}1{,}405$ lines of C++).
Table~\ref{tab:source_mods} provides the complete modification summary.



\begin{table}[!htbp]
\centering
\caption{Upstream CARLA source modifications.}
\label{tab:source_mods}
\small
\setlength{\tabcolsep}{3pt}
\renewcommand{\arraystretch}{1.08}
\resizebox{\columnwidth}{!}{%
\begin{tabular}{@{}lp{4.8cm}@{}}
\toprule
\textbf{File} & \textbf{Modification} \\
\midrule
\texttt{CarlaGameModeBase.h}
  & \texttt{friend} declaration for \texttt{CARLAAirGameMode} \\
\texttt{CarlaEpisode.h}
  & Sensor tagging visibility: \texttt{private} $\to$ \texttt{protected} \\
\texttt{CarlaGameModeBase.cpp}
  & State flag assignment (1 line) \\
\midrule
\multicolumn{2}{@{}l}{\textit{New integration layer (no upstream modification)}} \\
\texttt{CARLAAirGameMode}
  & Unified game mode class; ${\sim}1{,}405$ lines C++ \\
\bottomrule
\end{tabular}}%
\end{table}

\subsection{Custom Asset Import}
\label{app:assets}

CARLA-Air supports importing custom 3D assets (robot platforms, vehicles, UAV
models, and environment objects) through an asset pipeline built on Unreal
Engine's content framework. Imported assets are registered as spawnable actor
classes and become accessible through the standard CARLA Python API
(\texttt{world.spawn\_actor()}). Once registered, custom assets share the same
physics tick, rendering pass, and sensor visibility as all built-in actors,
ensuring full consistency within the joint simulation environment. 


\appendix
\section{Appendix}
\subsection{System Configuration}
\label{app:config}


\begin{figure}[!htbp]
\centering
\includegraphics[width=\columnwidth, height=0.3\textheight, keepaspectratio]{figures/asset.jpg}
\caption{%
  Custom assets imported into CARLA-Air through the extensible asset pipeline.
  \textbf{Top:} a four-wheeled mobile robot with onboard LiDAR, imported from
  an external FBX model.
  \textbf{Bottom:} a custom electric sport car with user-defined vehicle
  dynamics.
  Both assets operate within the shared simulation world alongside all
  built-in CARLA traffic and AirSim aerial agents, and are visible to all
  sensor modalities.%
}
\label{fig:custom_assets}
\end{figure}


\paragraph{Reference hardware.}
All experiments in this report were conducted on the following configuration:
Ubuntu 20.04/22.04~LTS, NVIDIA RTX~A4000 (16\,GB GDDR6), AMD Ryzen~7 5800X
(8-core, 4.7\,GHz), 32\,GB DDR4-3200.

\paragraph{Software stack.}
CARLA~0.9.16, AirSim~1.8.1 (final stable open-source release),
Unreal Engine~4.26, Python~3.8+.

\paragraph{Network configuration.}
Table~\ref{tab:ports} lists the default port assignments.
Both RPC servers bind to \texttt{localhost} by default; remote connections
require explicit IP configuration.

\begin{table}[!htbp]
\centering
\caption{Default network port assignments.}
\label{tab:ports}
\footnotesize
\setlength{\tabcolsep}{4pt}
\renewcommand{\arraystretch}{1.08}
\begin{tabular}{@{}llr@{}}
\toprule
\textbf{Service} & \textbf{Protocol} & \textbf{Port} \\
\midrule
CARLA RPC Server   & TCP & 2000  \\
CARLA Streaming    & UDP & 2001  \\
AirSim RPC Server  & TCP & 41451 \\
\bottomrule
\end{tabular}
\end{table}

\paragraph{Distribution.}
The prebuilt binary package is approximately 19\,GB and includes a one-command
launcher (\texttt{CarlaAir.sh}).
The source distribution is approximately 651\,MB and is released under the MIT
license.


\subsubsection{API Compatibility Summary}
\label{app:api}

Table~\ref{tab:api_compat} summarizes the API compatibility status and test
coverage of the current CARLA-Air release.
All 89 automated CARLA API tests pass without modification; the full AirSim
flight control and sensor access API has been verified through manual and
scripted testing.
A total of 63 ROS\,2 topics are published across both simulation backends.

\begin{table}[!htbp]
\centering
\caption{API compatibility and test coverage.}
\label{tab:api_compat}
\small
\setlength{\tabcolsep}{3pt}
\renewcommand{\arraystretch}{1.08}
\resizebox{\columnwidth}{!}{%
\begin{tabular}{@{}lp{4.5cm}@{}}
\toprule
\textbf{Component} & \textbf{Status} \\
\midrule
CARLA API           & 89/89 automated tests passing \\
AirSim API          & Full flight control and sensor access verified \\
ROS\,2 topics       & 63 total (43 CARLA + 14 AirSim + 6 generic) \\
\midrule
\multicolumn{2}{@{}l}{\textit{Key upstream scripts confirmed working}} \\
\texttt{manual\_control.py}    & CARLA manual driving interface \\
\texttt{automatic\_control.py} & CARLA autopilot demonstration \\
\texttt{dynamic\_weather.py}   & CARLA weather preset cycling \\
\texttt{hello\_drone.py}       & AirSim basic flight demonstration \\
\bottomrule
\end{tabular}}%
\end{table}


\subsection{Upstream Source Modifications}
\label{app:source_mods}

CARLA-Air is designed to minimize modifications to the upstream CARLA codebase.
The integration touches only two header files and one source file, totaling
approximately 35 lines of changes.
All other integration code is purely additive, residing within the aerial
simulation plugin as the \texttt{CARLAAirGameMode} class
(${\sim}1{,}405$ lines of C++).
Table~\ref{tab:source_mods} provides the complete modification summary.



\begin{table}[!htbp]
\centering
\caption{Upstream CARLA source modifications.}
\label{tab:source_mods}
\small
\setlength{\tabcolsep}{3pt}
\renewcommand{\arraystretch}{1.08}
\resizebox{\columnwidth}{!}{%
\begin{tabular}{@{}lp{4.8cm}@{}}
\toprule
\textbf{File} & \textbf{Modification} \\
\midrule
\texttt{CarlaGameModeBase.h}
  & \texttt{friend} declaration for \texttt{CARLAAirGameMode} \\
\texttt{CarlaEpisode.h}
  & Sensor tagging visibility: \texttt{private} $\to$ \texttt{protected} \\
\texttt{CarlaGameModeBase.cpp}
  & State flag assignment (1 line) \\
\midrule
\multicolumn{2}{@{}l}{\textit{New integration layer (no upstream modification)}} \\
\texttt{CARLAAirGameMode}
  & Unified game mode class; ${\sim}1{,}405$ lines C++ \\
\bottomrule
\end{tabular}}%
\end{table}

\subsection{Custom Asset Import}
\label{app:assets}

CARLA-Air supports importing custom 3D assets (robot platforms, vehicles, UAV
models, and environment objects) through an asset pipeline built on Unreal
Engine's content framework. Imported assets are registered as spawnable actor
classes and become accessible through the standard CARLA Python API
(\texttt{world.spawn\_actor()}). Once registered, custom assets share the same
physics tick, rendering pass, and sensor visibility as all built-in actors,
ensuring full consistency within the joint simulation environment. 


\appendix
\section{Appendix}
\subsection{System Configuration}
\label{app:config}


\begin{figure}[!htbp]
\centering
\includegraphics[width=\columnwidth, height=0.3\textheight, keepaspectratio]{figures/asset.jpg}
\caption{%
  Custom assets imported into CARLA-Air through the extensible asset pipeline.
  \textbf{Top:} a four-wheeled mobile robot with onboard LiDAR, imported from
  an external FBX model.
  \textbf{Bottom:} a custom electric sport car with user-defined vehicle
  dynamics.
  Both assets operate within the shared simulation world alongside all
  built-in CARLA traffic and AirSim aerial agents, and are visible to all
  sensor modalities.%
}
\label{fig:custom_assets}
\end{figure}


\paragraph{Reference hardware.}
All experiments in this report were conducted on the following configuration:
Ubuntu 20.04/22.04~LTS, NVIDIA RTX~A4000 (16\,GB GDDR6), AMD Ryzen~7 5800X
(8-core, 4.7\,GHz), 32\,GB DDR4-3200.

\paragraph{Software stack.}
CARLA~0.9.16, AirSim~1.8.1 (final stable open-source release),
Unreal Engine~4.26, Python~3.8+.

\paragraph{Network configuration.}
Table~\ref{tab:ports} lists the default port assignments.
Both RPC servers bind to \texttt{localhost} by default; remote connections
require explicit IP configuration.

\begin{table}[!htbp]
\centering
\caption{Default network port assignments.}
\label{tab:ports}
\footnotesize
\setlength{\tabcolsep}{4pt}
\renewcommand{\arraystretch}{1.08}
\begin{tabular}{@{}llr@{}}
\toprule
\textbf{Service} & \textbf{Protocol} & \textbf{Port} \\
\midrule
CARLA RPC Server   & TCP & 2000  \\
CARLA Streaming    & UDP & 2001  \\
AirSim RPC Server  & TCP & 41451 \\
\bottomrule
\end{tabular}
\end{table}

\paragraph{Distribution.}
The prebuilt binary package is approximately 19\,GB and includes a one-command
launcher (\texttt{CarlaAir.sh}).
The source distribution is approximately 651\,MB and is released under the MIT
license.


\subsubsection{API Compatibility Summary}
\label{app:api}

Table~\ref{tab:api_compat} summarizes the API compatibility status and test
coverage of the current CARLA-Air release.
All 89 automated CARLA API tests pass without modification; the full AirSim
flight control and sensor access API has been verified through manual and
scripted testing.
A total of 63 ROS\,2 topics are published across both simulation backends.

\begin{table}[!htbp]
\centering
\caption{API compatibility and test coverage.}
\label{tab:api_compat}
\small
\setlength{\tabcolsep}{3pt}
\renewcommand{\arraystretch}{1.08}
\resizebox{\columnwidth}{!}{%
\begin{tabular}{@{}lp{4.5cm}@{}}
\toprule
\textbf{Component} & \textbf{Status} \\
\midrule
CARLA API           & 89/89 automated tests passing \\
AirSim API          & Full flight control and sensor access verified \\
ROS\,2 topics       & 63 total (43 CARLA + 14 AirSim + 6 generic) \\
\midrule
\multicolumn{2}{@{}l}{\textit{Key upstream scripts confirmed working}} \\
\texttt{manual\_control.py}    & CARLA manual driving interface \\
\texttt{automatic\_control.py} & CARLA autopilot demonstration \\
\texttt{dynamic\_weather.py}   & CARLA weather preset cycling \\
\texttt{hello\_drone.py}       & AirSim basic flight demonstration \\
\bottomrule
\end{tabular}}%
\end{table}


\subsection{Upstream Source Modifications}
\label{app:source_mods}

CARLA-Air is designed to minimize modifications to the upstream CARLA codebase.
The integration touches only two header files and one source file, totaling
approximately 35 lines of changes.
All other integration code is purely additive, residing within the aerial
simulation plugin as the \texttt{CARLAAirGameMode} class
(${\sim}1{,}405$ lines of C++).
Table~\ref{tab:source_mods} provides the complete modification summary.



\begin{table}[!htbp]
\centering
\caption{Upstream CARLA source modifications.}
\label{tab:source_mods}
\small
\setlength{\tabcolsep}{3pt}
\renewcommand{\arraystretch}{1.08}
\resizebox{\columnwidth}{!}{%
\begin{tabular}{@{}lp{4.8cm}@{}}
\toprule
\textbf{File} & \textbf{Modification} \\
\midrule
\texttt{CarlaGameModeBase.h}
  & \texttt{friend} declaration for \texttt{CARLAAirGameMode} \\
\texttt{CarlaEpisode.h}
  & Sensor tagging visibility: \texttt{private} $\to$ \texttt{protected} \\
\texttt{CarlaGameModeBase.cpp}
  & State flag assignment (1 line) \\
\midrule
\multicolumn{2}{@{}l}{\textit{New integration layer (no upstream modification)}} \\
\texttt{CARLAAirGameMode}
  & Unified game mode class; ${\sim}1{,}405$ lines C++ \\
\bottomrule
\end{tabular}}%
\end{table}

\subsection{Custom Asset Import}
\label{app:assets}

CARLA-Air supports importing custom 3D assets (robot platforms, vehicles, UAV
models, and environment objects) through an asset pipeline built on Unreal
Engine's content framework. Imported assets are registered as spawnable actor
classes and become accessible through the standard CARLA Python API
(\texttt{world.spawn\_actor()}). Once registered, custom assets share the same
physics tick, rendering pass, and sensor visibility as all built-in actors,
ensuring full consistency within the joint simulation environment. 


\appendix
\section{Appendix}
\subsection{System Configuration}
\label{app:config}


\begin{figure}[!htbp]
\centering
\includegraphics[width=\columnwidth, height=0.3\textheight, keepaspectratio]{figures/asset.jpg}
\caption{%
  Custom assets imported into CARLA-Air through the extensible asset pipeline.
  \textbf{Top:} a four-wheeled mobile robot with onboard LiDAR, imported from
  an external FBX model.
  \textbf{Bottom:} a custom electric sport car with user-defined vehicle
  dynamics.
  Both assets operate within the shared simulation world alongside all
  built-in CARLA traffic and AirSim aerial agents, and are visible to all
  sensor modalities.%
}
\label{fig:custom_assets}
\end{figure}


\paragraph{Reference hardware.}
All experiments in this report were conducted on the following configuration:
Ubuntu 20.04/22.04~LTS, NVIDIA RTX~A4000 (16\,GB GDDR6), AMD Ryzen~7 5800X
(8-core, 4.7\,GHz), 32\,GB DDR4-3200.

\paragraph{Software stack.}
CARLA~0.9.16, AirSim~1.8.1 (final stable open-source release),
Unreal Engine~4.26, Python~3.8+.

\paragraph{Network configuration.}
Table~\ref{tab:ports} lists the default port assignments.
Both RPC servers bind to \texttt{localhost} by default; remote connections
require explicit IP configuration.

\begin{table}[!htbp]
\centering
\caption{Default network port assignments.}
\label{tab:ports}
\footnotesize
\setlength{\tabcolsep}{4pt}
\renewcommand{\arraystretch}{1.08}
\begin{tabular}{@{}llr@{}}
\toprule
\textbf{Service} & \textbf{Protocol} & \textbf{Port} \\
\midrule
CARLA RPC Server   & TCP & 2000  \\
CARLA Streaming    & UDP & 2001  \\
AirSim RPC Server  & TCP & 41451 \\
\bottomrule
\end{tabular}
\end{table}

\paragraph{Distribution.}
The prebuilt binary package is approximately 19\,GB and includes a one-command
launcher (\texttt{CarlaAir.sh}).
The source distribution is approximately 651\,MB and is released under the MIT
license.


\subsubsection{API Compatibility Summary}
\label{app:api}

Table~\ref{tab:api_compat} summarizes the API compatibility status and test
coverage of the current CARLA-Air release.
All 89 automated CARLA API tests pass without modification; the full AirSim
flight control and sensor access API has been verified through manual and
scripted testing.
A total of 63 ROS\,2 topics are published across both simulation backends.

\begin{table}[!htbp]
\centering
\caption{API compatibility and test coverage.}
\label{tab:api_compat}
\small
\setlength{\tabcolsep}{3pt}
\renewcommand{\arraystretch}{1.08}
\resizebox{\columnwidth}{!}{%
\begin{tabular}{@{}lp{4.5cm}@{}}
\toprule
\textbf{Component} & \textbf{Status} \\
\midrule
CARLA API           & 89/89 automated tests passing \\
AirSim API          & Full flight control and sensor access verified \\
ROS\,2 topics       & 63 total (43 CARLA + 14 AirSim + 6 generic) \\
\midrule
\multicolumn{2}{@{}l}{\textit{Key upstream scripts confirmed working}} \\
\texttt{manual\_control.py}    & CARLA manual driving interface \\
\texttt{automatic\_control.py} & CARLA autopilot demonstration \\
\texttt{dynamic\_weather.py}   & CARLA weather preset cycling \\
\texttt{hello\_drone.py}       & AirSim basic flight demonstration \\
\bottomrule
\end{tabular}}%
\end{table}


\subsection{Upstream Source Modifications}
\label{app:source_mods}

CARLA-Air is designed to minimize modifications to the upstream CARLA codebase.
The integration touches only two header files and one source file, totaling
approximately 35 lines of changes.
All other integration code is purely additive, residing within the aerial
simulation plugin as the \texttt{CARLAAirGameMode} class
(${\sim}1{,}405$ lines of C++).
Table~\ref{tab:source_mods} provides the complete modification summary.



\begin{table}[!htbp]
\centering
\caption{Upstream CARLA source modifications.}
\label{tab:source_mods}
\small
\setlength{\tabcolsep}{3pt}
\renewcommand{\arraystretch}{1.08}
\resizebox{\columnwidth}{!}{%
\begin{tabular}{@{}lp{4.8cm}@{}}
\toprule
\textbf{File} & \textbf{Modification} \\
\midrule
\texttt{CarlaGameModeBase.h}
  & \texttt{friend} declaration for \texttt{CARLAAirGameMode} \\
\texttt{CarlaEpisode.h}
  & Sensor tagging visibility: \texttt{private} $\to$ \texttt{protected} \\
\texttt{CarlaGameModeBase.cpp}
  & State flag assignment (1 line) \\
\midrule
\multicolumn{2}{@{}l}{\textit{New integration layer (no upstream modification)}} \\
\texttt{CARLAAirGameMode}
  & Unified game mode class; ${\sim}1{,}405$ lines C++ \\
\bottomrule
\end{tabular}}%
\end{table}

\subsection{Custom Asset Import}
\label{app:assets}

CARLA-Air supports importing custom 3D assets (robot platforms, vehicles, UAV
models, and environment objects) through an asset pipeline built on Unreal
Engine's content framework. Imported assets are registered as spawnable actor
classes and become accessible through the standard CARLA Python API
(\texttt{world.spawn\_actor()}). Once registered, custom assets share the same
physics tick, rendering pass, and sensor visibility as all built-in actors,
ensuring full consistency within the joint simulation environment. 
